\chapter{Conclusion}
\label{ch:conclusion}

\section{Objectives revisited}

Table~\ref{tab:objectives} maps each objective from Section~\ref{sec:objectives}
to the evidence bearing on it. Five were met; one was met in a form different from
the one proposed.

\begin{table}[htbp]
\centering
\footnotesize
\caption[Objectives against evidence]{Objectives against the evidence in this
report.}
\label{tab:objectives}
\begin{tabular}{@{}clp{5.5cm}l@{}}
\toprule
& \textbf{Objective} & \textbf{Evidence} & \textbf{Status} \\
\midrule
O1 & Three agents + summariser & \S\ref{sec:grounding}, Table~\ref{tab:agent-config};
     interests traced to Doc 1 categories & Met \\
O2 & Deterministic orchestration, validated output & \S\ref{sec:orchestration};
     structural compliance 0.9986 over 2{,}865 responses & Met \\
O3 & Conditionally verifying court agent & \S\ref{sec:court-design};
     0 of 9 flagged instances Court-originated & Met \\
O4 & Real documents to structured scenarios & \S\ref{sec:extraction};
     23 cases ingested through the user upload path & Met \\
O5 & Verified corpus, baselines, ablations & \S\ref{sec:baselines}--\S\ref{sec:ablation};
     13 of 14 merits dispositions matched & Met \\
O6 & Something the client can use & \S\ref{sec:risk-screen};
     rule-based screen, not the simulator & Met, reframed \\
\bottomrule
\end{tabular}
\end{table}

\section{What the research questions turned out to have as answers}

\textbf{RQ1, on prompt engineering without fine-tuning.} Yes, with a caveat that
only appeared under testing. Agents specified purely through role prompts and
injected context do sustain distinct, role-appropriate positions: the authority
concedes rarely on opening and about half the time thereafter, the bidder concedes
far less, and the two report opposite outcomes against their own declared
fallbacks. The caveat is that a prompt developed against one dispute shape does
not generalise to another. Confronted with a case where nobody was ever scored,
both negotiating agents invented a scoring dispute, because that was the only
script they had.

\textbf{RQ2, on which part of the protocol is load-bearing.} The adjudicator, and
nothing else. Removing it deadlocks every run. That result is close to trivial in
retrospect --- nothing else in the graph can set the resolution flag --- but it is
the only unambiguous architectural finding in this project, and it is worth more
than the accuracy figures, because the accuracy figures do not distinguish this
system from a single prompt.

\textbf{RQ3, on constraining an LLM adjudicator's arithmetic.} Largely yes, and
the boundary is now mapped. Gating exact computation on the presence of both a
formula and every input it needs, and banning unstated numbers even when hedged as
illustrative, holds across 425 qualitative assessments with no Court-originated
fabrication among the nine flagged. It does not hold universally: a documented
instance exists of the Court inventing an entire illustrative dataset on a case it
misclassified as numeric, and the arithmetic defect found on the client's own
scoring table resisted three prompt revisions. The discipline constrains what the
adjudicator computes; it does not constrain what premise it computes on.

\textbf{RQ4, on evaluating in the absence of ground truth.} The answer this
project arrived at is a stack of partial checks with explicit sample sizes, and a
refusal to combine them into a single score. The stack has to include an
architecture check, because a system can agree with reality for reasons that have
nothing to do with its design; it has to include fabrication screens, because a
correct conclusion reached from invented evidence looks identical from outside to a
correct conclusion; and it has to report negative results, because the two most
informative findings here --- the zero-shot tie and the readability failure ---
both count against the artefact.

\section{Reflection on the approach}

The scope was wide: five agent types, a state-machine orchestrator, a document
ingestion path, a streaming API, a client-facing web application, a rule-based
risk screen, a 23-case verified legal corpus, and eight analysis scripts, held
together by 92 tests. The complexity that actually mattered, though, was not any
of those individually. It was that the domain punishes exactly the thing
generative models are best at. A language model asked about a procurement dispute
will produce something fluent, structured and legally shaped, and the harder the
case, the more confident it sounds. Almost every engineering decision in
Chapter~\ref{ch:methodology} is downstream of that: schema-validated output so
malformed responses cannot pass silently, per-response compliance counters so
repairs are visible, an extraction agent instructed as firmly against dropping
real numbers as against inventing them, a minimum source-length guard because an
empty PDF produced a complete fictional dispute, and a conditional verification
protocol whose entire purpose is to make the adjudicator say ``I cannot compute
this'' when it cannot.

Two things I would do differently. First, I would build the corpus before building
the agents. The synthetic scenarios written early on were AI-authored, which made
every result derived from them circular, and they were eventually deleted along
with seven directories of evaluation output --- a month of work that established
nothing. Second, I would have tested on documents I did not choose much sooner.
Two of the most informative defects in this project came from an hour spent
running a Fusion21 employee's own files, and both were in code paths that months
of my own test cases had never exercised, because I had unknowingly been selecting
inputs shaped like the failures I already understood.

On execution quality, the honest position is mixed and the mix is informative. The
structural layer is solid: 0.9986 clean structured responses, a green test suite,
every result reproducible from committed logs, and provenance carried through to
the analytics interface so a partial batch cannot be mistaken for a finished one.
The reasoning layer is where the limits are, and they are documented rather than
smoothed over --- a 23.1\% wrong-regime citation rate, a summariser that fails its
own accessibility claim, an outcome vocabulary that cannot express three of the
real remedies in its own corpus, and one arithmetic defect that prompt engineering
did not close.

\section{Future work}

Four directions follow directly from findings above rather than from speculation.

\textbf{Ground the citations.} A 23.1\% wrong-regime rate is the strongest
argument in this report for the retrieval layer that was scoped out. Retrieval
over the Procurement Act 2023 and the TCC guidance targets exactly the failure
observed: the model does not know which regime governs a given year, and no
instruction can supply that knowledge, whereas retrieval can.

\textbf{Check premises the way arithmetic is checked.} The Court agent verifies
numbers against the scenario record but accepts factual assertions from the
negotiating agents without checking them against the same record. The Faraday
result shows what that costs. A grounding check applied to negotiating-agent
statements, using the machinery the numeric screen already provides, is a
concrete and well-evidenced next step.

\textbf{Extend the remedy vocabulary beyond scoring challenges.} The corpus itself
contains three real remedies the system cannot express: a declaration of
ineffectiveness, financial penalties on the authority, and a court-determined
re-ranking. Each is a documented gap on a named case, not a hypothetical.

\textbf{Calibrate the risk screen against real outcomes.} This is the only ask
this project has of Fusion21, and it is the highest-value item on this list.
Anonymised pre-action data from their member base --- the same profile fields the
screen already requests, plus whether a challenge followed --- would let a
severity-weighted heuristic be replaced or validated by a fitted model, and would
let the risk score be reported as a probability instead of an ordering. It is also
the only component here whose central claim is directly falsifiable, which is
precisely why it deserves the data.

The broader conclusion is narrower than the one I expected to write. A
multi-agent LLM system can conduct a procedurally coherent procurement dispute
negotiation, and can be constrained not to invent the arithmetic it reasons from.
It cannot yet be shown to reason better than a single well-posed prompt, and it
can still be led into confident nonsense by a premise it was handed rather than one
it computed. The value in this work is less the simulator than the evidence about
what such a system does when nobody is checking --- and the demonstration that
finding out requires screens built specifically to catch a system that is wrong in
a fluent and plausible way.
